\section{Background and Architectural Comparison}
\label{sec:background}

This section examines the architectural foundations of Kubernetes-native chaos engineering platforms. We first contrast the control-plane and node-level execution models of Chaos Mesh and LitmusChaos across their lifecycle stages, and subsequently review related comparative studies to establish the empirical research gap addressed by this work.

\subsection{Architectural Divergence: Chaos Mesh vs.\ LitmusChaos}
Although both Chaos Mesh and LitmusChaos are CNCF Incubating projects designed around the Kubernetes Operator pattern~\cite{owotogbe2026tenyears, onkalo2025improving}, their internal control and data planes exhibit distinct architectural trade-offs, illustrated in Figure~\ref{fig:cm_vs_litmus_arch} and summarized in Table~\ref{tab:architectural_comparison}.

\begin{figure}[!htbp]
\centering
\begin{tikzpicture}[
    node distance=0.36cm and 0.25cm,
    font=\sffamily\scriptsize,
    basebox/.style={draw, rounded corners=2pt, align=center, text width=3.6cm, minimum width=3.75cm, minimum height=0.55cm, inner sep=2.5pt},
    cmbox/.style={basebox, fill=blue!8, draw=blue!65!black, thick},
    lcbox/.style={basebox, fill=orange!10, draw=orange!75!black, thick},
    targetbox/.style={basebox, fill=teal!10, draw=teal!70!black},
    arrow/.style={-{Stealth[scale=0.75]}, thick, draw=black!75}
]
    % Chaos Mesh Side
    \node[cmbox] (cm_crd) {\textbf{Chaos Mesh CRD}\\(\texttt{Network/Pod/StressChaos})};
    \node[cmbox, below=0.35cm of cm_crd] (cm_ctrl) {\texttt{chaos-controller-manager}\\(Admission \& Reconciliation)};
    \node[cmbox, below=0.35cm of cm_ctrl] (cm_daemon) {\texttt{chaos-daemon} (Resident DaemonSet)\\Direct gRPC $\rightarrow$ \texttt{nsenter} / \texttt{cgroups}};
    \node[targetbox, below=0.35cm of cm_daemon] (cm_pod) {\textbf{Target Pod (\texttt{otel-demo})}\\Immediate Kernel Perturbation};

    % LitmusChaos Side
    \node[lcbox, right=0.32cm of cm_crd] (lc_crd) {\textbf{LitmusChaos CRDs}\\(\texttt{ChaosEngine} + \texttt{ChaosExperiment})};
    \node[lcbox, below=0.35cm of lc_crd] (lc_op) {\texttt{chaos-operator} (+ Optional \texttt{ChaosCenter})\\CRD Reconciliation \& Runner Spec};
    \node[lcbox, below=0.35cm of lc_op] (lc_runner) {Spawns Ephemeral Pods per Trial:\\\texttt{chaos-runner} $\rightarrow$ \texttt{go-runner} Job};
    \node[targetbox, below=0.35cm of lc_runner] (lc_pod) {\textbf{Target Pod (\texttt{otel-demo})}\\Fault Injection + \texttt{Litmus Probes}};

    \draw[arrow, blue!70!black] (cm_crd) -- (cm_ctrl);
    \draw[arrow, blue!70!black] (cm_ctrl) -- node[right, font=\sffamily\tiny] {gRPC} (cm_daemon);
    \draw[arrow, blue!70!black] (cm_daemon) -- node[right, font=\sffamily\tiny] {$\Delta T_{\text{actuate}} < 1\,\text{s}$} (cm_pod);

    \draw[arrow, orange!80!black] (lc_crd) -- (lc_op);
    \draw[arrow, orange!80!black] (lc_op) -- node[right, font=\sffamily\tiny] {Pod Create} (lc_runner);
    \draw[arrow, orange!80!black] (lc_runner) -- node[right, font=\sffamily\tiny] {$\Delta T_{\text{actuate}} \approx 11.8\,\text{s}$} (lc_pod);
\end{tikzpicture}
\caption{Execution lifecycle comparison: Chaos Mesh's resident DaemonSet injection vs.\ LitmusChaos's ephemeral runner pod orchestration.}
\Description{Side-by-side diagram comparing Chaos Mesh direct DaemonSet injection and LitmusChaos ephemeral runner pod orchestration.}
\label{fig:cm_vs_litmus_arch}
\end{figure}

\begin{table*}[!tbp]
\centering
\scriptsize
\setlength{\tabcolsep}{4pt}
\renewcommand{\arraystretch}{1.15}
\caption{Architectural and Functional Comparison Between Chaos Mesh and LitmusChaos}
\label{tab:architectural_comparison}
\begin{tabularx}{\textwidth}{p{2.65cm} X X}
\toprule
\textbf{Architectural Dimension} & \textbf{Chaos Mesh (v2.6+ / v2.8)} & \textbf{LitmusChaos (v3.x)} \\
\midrule
\textbf{Core Control Plane} & \texttt{chaos-controller-manager}, \texttt{chaos-dashboard}, \texttt{chaos-dns-server} (Optional SQLite/MySQL persistence). & Standalone \texttt{chaos-operator} (\texttt{litmus-core}) or full \texttt{ChaosCenter} (Frontend + Auth + GraphQL + \texttt{MongoDB} + \texttt{chaos-exporter}). \\
\textbf{Node Execution Mechanism} & Pre-warmed, node-resident \texttt{chaos-daemon} DaemonSet executing direct gRPC-triggered \texttt{nsenter} syscalls. & On-demand ephemeral \texttt{chaos-runner} pod spawning transient \texttt{go-runner} Kubernetes Jobs per experiment ($+2$ pods/trial). \\
\textbf{Fault Catalog Scope} & 11 core fault categories (Pod, Network, Stress, IO, DNS, Time/Clock Skew, Kernel/Slab, HTTP, JVM, Block, Physical Node via \texttt{chaosd}) \cite{onkalo2025improving}. & 53+ ChaosHub experiments spanning Kubernetes (Pod, Node, Network, Stress, DNS) and Multi-Cloud IaaS (AWS EC2/EBS/RDS, Azure, GCP, VMware) \cite{onkalo2025improving, opara2025chaos}. \\
\textbf{Native Hypothesis Verification} & \texttt{StatusCheck} CRD supporting synchronous/continuous HTTP status/body polling; lacks composite resilience scoring. & Declarative \textbf{Litmus Probes} (\texttt{httpProbe}, \texttt{cmdProbe}, \texttt{k8sProbe}, \texttt{promProbe}) in 4 modes (\texttt{SOT}, \texttt{EOT}, \texttt{Edge}, \texttt{Continuous}) + quantitative \textbf{Resilience Score (\%)} \cite{onkalo2025improving}. \\
\textbf{Observability Export} & Prometheus metrics via controller/daemon endpoints; Chaos Mesh Dashboard events. & Dedicated \texttt{chaos-exporter} exposing per-experiment Prometheus metrics and OpenTelemetry hooks. \\
\bottomrule
\end{tabularx}
\end{table*}

\textbf{Chaos Mesh Architecture:} Chaos Mesh prioritizes low-latency, kernel-proximate fault injection within Kubernetes clusters. Its control plane (\texttt{chaos-controller-manager}) watches type-specific CRDs (\texttt{NetworkChaos}, \texttt{PodChaos}, \texttt{StressChaos}, \texttt{IOChaos}, \texttt{TimeChaos}, \texttt{KernelChaos}) via admission webhooks. Because the privileged \texttt{chaos-daemon} DaemonSet is already running on every worker node, injecting a fault requires zero new pod scheduling: \texttt{chaos-controller-manager} invokes a gRPC call to the target node's \texttt{chaos-daemon}, which immediately attaches to the target container's network, PID, or mount namespace~\cite{onkalo2025improving, owotogbe2026tenyears}.

\textbf{LitmusChaos Architecture:} LitmusChaos is architected as a comprehensive enterprise resilience management platform~\cite{onkalo2025improving, redkar2026enhancing, gurudu2026chaos}. Its control plane (\texttt{ChaosCenter}) relies on a persistent \texttt{MongoDB} instance to store multi-tenant chaos workflows, ChaosHub templates, and historical Resilience Scores. When a \texttt{ChaosEngine} CRD binds a target application to a \texttt{ChaosExperiment} (e.g., \texttt{pod-network-latency}, \texttt{pod-delete}, or \texttt{pod-cpu-hog}), the \texttt{chaos-operator} creates a \texttt{chaos-runner} pod in the target namespace. The \texttt{chaos-runner} reads the experiment parameters and creates a secondary \texttt{go-runner} helper pod scheduled onto the target microservice's node to execute the container/network perturbation and evaluate attached \textbf{Litmus Probes}~\cite{onkalo2025improving}. Upon completion, a \texttt{ChaosResult} CRD is updated and exported via \texttt{chaos-exporter}.

\subsection{Related Comparative Literature}
While several studies employ either LitmusChaos~\cite{gurudu2026chaos, redkar2026enhancing} or Chaos Mesh~\cite{opara2025chaos, nigam2026aiaugmented} in isolation, comparative analyses have remained strictly qualitative. Mailewa et al.~\cite{mailewa2025review} reviewed six chaos tools (including Chaos Mesh and Litmus) conceptually, while Opara et al.~\cite{opara2025chaos} and Ambani~\cite{ambani2025systematic} compared their feature matrices across multi-cloud reach and CI/CD integration. Most recently, Onkalo~\cite{onkalo2025improving} conducted a dedicated feature comparison between LitmusChaos v3.22.0 and Chaos Mesh v2.8.0 against ISO/IEC 25010 reliability criteria, concluding that LitmusChaos provides broader cloud-provider coverage ($53$ experiments) and superior native probes/CI-CD REST APIs, whereas Chaos Mesh provides deeper kernel/time injection and lower setup complexity. However, Onkalo~\cite{onkalo2025improving} performed no empirical experiments, explicitly calling for quantitative benchmarking of both tools under realistic workloads. This study directly fulfills that research need.